\chapter{Podstawy teoretyczne i technologiczne}\label{ch:podstawy-technologiczne}

Projektowane narzędzie operuje na styku kilku podsystemów jądra oraz przestrzeni użytkownika systemu Microsoft Windows, a w części dotyczącej kopii zapasowych wykorzystuje kryptograficzne protokoły komunikacyjne systemu Linux. Niniejszy rozdział przedstawia fundamenty teoretyczne i technologiczne, które bezpośrednio determinują decyzje architektoniczne opisane w rozdziale~\ref{ch:architektura-i-dzialanie}.

\section{Architektura Zapory systemu Windows i platformy WFP}\label{sec:architektura-zapory}

\subsection{Platforma filtrująca Windows Filtering Platform (WFP)}\label{subsec:architektura-wfp}

Począwszy od systemów Windows Vista oraz Windows Server 2008, architektura filtracji ruchu sieciowego została oparta na platformie WFP (\textit{Windows Filtering Platform}) \cite{ms-wfp}. W odróżnieniu od starszych filtrów pakietowych opartych na sterownikach pośrednich NDIS, WFP udostępnia wielowarstwowy szkielet przetwarzania pakietów zintegrowany bezpośrednio ze stosem sieciowym TCP/IP (\file{tcpip.sys}).

Architektura zapory dzieli się na dwie współpracujące ze sobą warstwy:
\begin{itemize}
    \item \textbf{Warstwa trybu użytkownika (\textit{User Mode})} - obejmuje usługę systemową \code{MpsSvc} (\textit{Windows Defender Firewall}) oraz podstawowy silnik filtrujący BFE (\textit{Base Filtering Engine}), działające w kontekście procesów \file{svchost.exe}. Usługa \code{MpsSvc} odpowiada za odczyt trwałych magazynów reguł, uwierzytelnianie IPsec, obsługę profili sieciowych oraz tłumaczenie wysokopoziomowych reguł zapory na niskopoziomowe filtry przekazywane do BFE.
    \item \textbf{Warstwa trybu jądra (\textit{Kernel Mode})} - obejmuje silnik filtrujący jądra zintegrowany ze sterownikiem \file{tcpip.sys} oraz sterownikiem lekkiego filtru NDIS (\file{wfplwfs.sys}). W warstwie tej znajdują się tzw.\ \textit{shims} (punkty zaczepienia na poszczególnych etapach przetwarzania stosu sieciowego, m.in.\ przy odbiorze ramki, routingu IP, nawiązywaniu połączenia warstwy transportowej \code{ALE_AUTH_CONNECT} oraz akceptacji gniazda \code{ALE_AUTH_RECV_ACCEPT}), które przekazują metadane pakietu do klasyfikatora reguł.
\end{itemize}

Klasyfikator jądra ocenia ruch sieciowy w oparciu o zestaw warunków dopasowania (adresy źródłowe i docelowe IPv4/IPv6, numery portów lokalnych i zdalnych, numer protokołu IP, ścieżkę pliku wykonywalnego aplikacji oraz identyfikator usługi systemowej). Domyślna polityka zapory Windows we wszystkich profilach zakłada \textbf{blokowanie ruchu przychodzącego} (\textit{Inbound}), który nie pasuje do żadnej jawnej reguły zezwalającej, oraz \textbf{zezwalanie na ruch wychodzący} (\textit{Outbound}), o ile nie zostanie on dopasowany do reguły blokującej \cite{ms-defender-firewall}. W przypadku konfliktu reguł w obrębie tego samego magazynu polityki zasada precedencji przyznaje wyższy priorytet akcji \code{Block} nad akcją \code{Allow} (z wyjątkiem reguł posiadających jawny atrybut nadrzędności uwierzytelnionej \textit{OverrideBlockRules}).

\subsection{Magazyn polityki w rejestrze systemowym i obiekty GPO}\label{subsec:struktura-regul}

Z punktu widzenia projektowania podsystemu audytu oraz kopii zapasowych kluczowe znaczenie ma fizyczne miejsce składowania konfiguracji zapory. Usługa \code{MpsSvc} przechowuje trwały stan polityki lokalnej w gałęzi rejestru systemowego \code{HKEY_LOCAL_MACHINE} (HKLM) pod kluczem usługowym \file{SYSTEM\CurrentControlSet\Services\SharedAccess\Parameters\FirewallPolicy}, natomiast reguły i ustawienia narzucone centralnie przez mechanizm zasad grupy (\textit{Group Policy Objects} - GPO) są odwzorowywane w drzewie zasad \file{SOFTWARE\Policies\Microsoft\WindowsFirewall}. Strukturę kluczowych podkluczy magazynu polityki zapory wykorzystywanych przez silnik audytowy przedstawiono w tabeli~\ref{tab:rejestr-zapory}.

\begin{table}[H]
    \centering
    \caption{Struktura magazynu polityki Zapory systemu Windows w rejestrze systemowym}\label{tab:rejestr-zapory}
    \tablefont
    \begin{tabularx}{\linewidth}{L{3.8cm} L{2.9cm} Y}
        \toprule
        \hcell{Podklucz rejestru} & \hcell{Prefiks audytu} & \hcell{Przechowywane obiekty konfiguracji} \\
        \midrule
        \code{FirewallRules} & \code{rule:} / \code{gporule:} & Reguły filtracji pakietów przychodzących i wychodzących; nazwa wartości to unikalny identyfikator reguły (\code{Name}), a dane to łańcuch w formacie potokowym. \\
        \code{ConSecRules} & \code{consec:} / \code{gpoconsec:} & Reguły zabezpieczeń połączeń (\textit{Connection Security Rules}) definiujące polityki uwierzytelniania i szyfrowania IPsec między hostami. \\
        \code{DomainProfile} & \code{profile:Domain} & Parametry globalne profilu domenowego (włączenie zapory \code{EnableFirewall}, domyślne akcje, powiadomienia) wraz z podkluczem \code{Logging}. \\
        \code{StandardProfile} & \code{profile:Private} & Parametry globalne profilu sieci prywatnej (historyczna nazwa klucza \code{StandardProfile} pochodzi z systemu Windows XP/Server 2003). \\
        \code{PublicProfile} & \code{profile:Public} & Parametry globalne profilu sieci publicznej wraz z ustawieniami dziennika pakietów odrzuconych. \\
        Pozostałe podklucze (\code{Mdm}, \code{RestrictedServices} i in.) & \code{setting:} / \code{gposetting:} & Dodatkowe reguły usług ograniczonych, konfiguracja zarządzania mobilnego MDM, dynamiczne słowa kluczowe oraz reguły wirtualizacji Hyper-V. \\
        \bottomrule
    \end{tabularx}
\end{table}

Każda reguła w kluczu \code{FirewallRules} jest zapisywana jako pojedyncza wartość tekstowa (\code{REG_SZ}) składająca się z nagłówka wersji schematu (np.\ \code{v2.33}) oraz par \code{Klucz=Wartość} rozdzielanych znakiem pionowej kreski (\code{|}), przykładowo:
\begin{center}
    {\footnotesize\code{v2.33|Action=Block|Active=TRUE|Dir=In|Protocol=6|Profile=Public|LPort=22|}}\\[-2pt]
    {\footnotesize\code{RA4=64.89.160.111|Name=01-OSO-Block-Inbound-222442|}}
\end{center}
Wewnątrz tego zapisu numery protokołów są reprezentowane zgodnie z rejestrem IANA (\code{6} dla TCP, \code{17} dla UDP, \code{1} dla ICMPv4, \code{58} dla ICMPv6, \code{256} dla dowolnego protokołu), a wielokrotne profile lub adresy IP występują jako powtarzające się tokeny (np.\ \code{Profile=Private|Profile=Public|} lub \code{RA4=...} dla IPv4 i \code{RA6=...} dla IPv6). Z kolei predefiniowane reguły systemowe używają pośrednich odwołań do zasobów bibliotek DLL w polach nazwy i opisu (np.\ \code{@FirewallAPI.dll,-32765}), które system rozwija do czytelnych napisów za pomocą funkcji natywnej \code{SHLoadIndirectString} z biblioteki \file{shlwapi.dll}.

\section{Profile sieciowe i usługa Network Location Awareness (NLA)}\label{sec:profile-sieciowe}

\subsection{Charakterystyka profili: domenowy, prywatny i publiczny}\label{subsec:charakterystyka-profili}

Aby umożliwić dynamiczne dostosowanie poziomu restrykcyjności zapory do aktualnego otoczenia sieciowego bez konieczności ręcznego przepisywania reguł, system Windows definiuje trzy niezależne profile bezpieczeństwa \cite{ms-defender-firewall}:
\begin{enumerate}
    \item \textbf{Profil domenowy (\textit{Domain})} - aktywowany wyłącznie wtedy, gdy interfejs sieciowy stacji roboczej posiada łączność z kontrolerem domeny usługi Active Directory, do której należy komputer, i pomyślnie przeprowadzi uwierzytelnienie domenowe. Jest to profil o najwyższym poziomie zaufania, przeznaczony dla wewnętrznych sieci korporacyjnych.
    \item \textbf{Profil prywatny (\textit{Private})} - przypisywany przez administratora zaufanym sieciom lokalnym (domowym lub małym sieciom biurowym), w których urządzenia są znane, a sieć jest odseparowana od Internetu za pomocą routera brzegowego z translacją adresów NAT. W profilu tym dozwolone jest m.in.\ wykrywanie sieci oraz udostępnianie plików.
    \item \textbf{Profil publiczny (\textit{Public})} - domyślny, najbardziej restrykcyjny profil stosowany dla wszystkich nowo napotkanych i niezidentyfikowanych sieci (kawiarnie, lotniska, hotspoty mobilne, bezpośrednie łącza publiczne). W profilu tym zapora powinna blokować wszelkie przychodzące próby nawiązania połączeń do usług lokalnych.
\end{enumerate}

\subsection{Mechanizm rozpoznawania lokalizacji sieciowej (NLA)}\label{subsec:mechanizm-nla}

Za automatyczne wyznaczanie profilu dla każdego aktywnego adaptera sieciowego odpowiada usługa systemowa \code{NlaSvc} (\textit{Network Location Awareness}) współpracująca z usługą listy sieci (\code{netprofm}). W momencie podłączenia kabla Ethernet lub skojarzenia z siecią Wi-Fi usługa NLA zbiera cechy charakterystyczne sieci (m.in.\ sufiks DNS połączenia, adres IP i adres fizyczny MAC bramy domyślnej, dostępność kontrolera domeny przez protokół LDAP/Kerberos), tworzy jej podpis kryptograficzny i porównuje go z zapamiętanymi profilami w rejestrze \code{HKLM\SOFTWARE\Microsoft\Windows NT\CurrentVersion\NetworkList\Profiles}.

Na tej podstawie usługa NLA przypisuje interfejsowi kategorię (\code{NetworkCategory}), a silnik BFE aktywuje wyłącznie reguły powiązane z danym profilem (lub wartością \code{Any}) \cite{ms-wfp}. Warto podkreślić, że w systemach z wieloma kartami sieciowymi profile są stosowane niezależnie dla każdego interfejsu (\textit{per-interface profile binding}): jeżeli komputer jest podłączony jednocześnie do prywatnej sieci LAN oraz publicznego modemu bezprzewodowego, ruch przychodzący przez interfejs publiczny podlega wyłącznie regułom przypisanym do profilu \code{Public}.

\section{Automatyzacja administracji zaporą w środowisku PowerShell}\label{sec:automatyzacja-powershell}

\subsection{Moduł NetSecurity, narzędzie netsh advfirewall oraz bezpośredni odczyt rejestru}\label{subsec:modul-netsecurity}

W środowisku Windows PowerShell 5.1 (\textit{Desktop Edition}) zarządzanie zaporą może odbywać się za pomocą trzech komplementarnych interfejsów technologicznych, z których każdy charakteryzuje się odmiennymi właściwościami wydajnościowymi i funkcjonalnymi \cite{ms-netsecurity-module, ms-netsh-advfirewall}:

\begin{enumerate}
    \item \textbf{Moduł \code{NetSecurity} (WMI/CIM)} - zestaw poleceń cmdlet przestrzeni WMI/CIM (m.in.\ \cmd{New-NetFirewallRule}, \cmd{Get-NetFirewallRule}) \cite{ms-new-netfirewallrule}. Parametry adresowe i portowe są rozbite na filtry obiektowe (\cmd{AddressFilter} oraz \cmd{PortFilter}), natomiast profil NLA odczytuje polecenie \cmd{Get-NetConnectionProfile}.
    \item \textbf{Narzędzie konsolowe \cmd{netsh advfirewall}} - natywny program komunikujący się z API zapory (\code{INetFwPolicy2}), będący jedynym wbudowanym narzędziem oferującym atomowy, binarny eksport oraz import całej polityki zapory (wraz z profilami i regułami IPsec) do pliku \file{.wfw} (\cmd{netsh advfirewall export} / \cmd{import}) \cite{ms-netsh-advfirewall}.
    \item \textbf{Bezpośredni odczyt magazynu rejestru (\code{Microsoft.Win32.Registry})} - pobranie około 500 reguł systemowych przez \cmd{Get-NetFirewallRule} wraz z filtrami adresów i portów wymaga ponad tysiąca zapytań CIM i trwa kilkanaście sekund. Bezpośrednie odczytanie i sparsowanie surowych wartości z klucza \code{FirewallRules} za pomocą klas .NET zajmuje zaledwie kilkadziesiąt milisekund (około 30-50~ms). Całościowy proces wyznaczenia różnic (\textit{diff}), parsowania obiektów GPO i atrybucji sprawcy zamyka się poniżej 1-2~s, realizując wymaganie \code{WNF-04}.
\end{enumerate}

W opracowanym systemie połączono zalety wszystkich trzech podejść: moduł \code{NetSecurity} służy do walidowanego tworzenia, filtrowania i usuwania reguł projektowych \code{OSO-*}, bezpośredni odczyt rejestru napędza szybkie migawki silnika audytowego, natomiast \cmd{netsh advfirewall} odpowiada za pełne zrzuty i odtwarzanie kopii zapasowych \file{.wfw}.

\subsection{Uprawnienia administratora, polityki wykonywania i listy kontroli dostępu (ACL)}\label{subsec:uprawnienia-executionpolicy}

Modyfikacja polityki zapory ogniowej oraz konfiguracja zadań harmonogramu wymagają podniesionego tokenu bezpieczeństwa procesu (\textit{elevated token}) uzyskiwanego przez mechanizm Kontroli Konta Użytkownika (UAC). W skryptach PowerShell wymóg ten jest egzekwowany na poziomie analizatora składniowego za pomocą dyrektywy preprocesora \code{\#Requires -RunAsAdministrator}, która uniemożliwia uruchomienie skryptu w nieuprzywilejowanej sesji. Równocześnie dyrektywy \code{\#Requires -Version 5.1} oraz \code{\#Requires -PSEdition Desktop} gwarantują dostępność pełnych klas .NET Framework wymaganych przez natywne interfejsy Win32.

Istotnym zagadnieniem z zakresu ochrony systemu operacyjnego jest zabezpieczenie plików skryptów oraz dzienników audytu przechowywanych w katalogu \file{C:\ProgramData}. Domyślna lista kontroli dostępu (DACL) katalogu \file{C:\ProgramData} zezwala zwykłym, nieuprzywilejowanym użytkownikom (grupa \code{Users}, identyfikator SID \code{S-1-5-32-545}) na tworzenie nowych podkatalogów i plików oraz ich modyfikację w ramach utworzonych przez siebie struktur. Gdyby podkatalog audytu lub kopia skryptu wykonywana przez zadanie systemowe odziedziczyły te uprawnienia lub zostały wcześniej utworzone przez zwykłego użytkownika (\textit{pre-creation attack}), lokalny użytkownik mógłby zmodyfikować plik bazowy \file{baseline.json}, usunąć dziennik audytu lub wstrzyknąć kod do skryptu uruchamianego z uprawnieniami \code{SYSTEM}.

Dlatego w projekcie zastosowano programową weryfikację właściciela oraz wyłączenie dziedziczenia uprawnień DACL (\code{SetAccessRuleProtection(\$true, \$false)}), przyznając pełną kontrolę (\code{FullControl}) wyłącznie dwóm zaufanym identyfikatorom SID:
\begin{itemize}
    \item \code{S-1-5-18} - konto lokalnego systemu operacyjnego (\code{NT AUTHORITY\SYSTEM}),
    \item \code{S-1-5-32-544} - wbudowana grupa administratorów (\code{BUILTIN\Administrators}).
\end{itemize}

\section{Integralność dzienników audytu i architektura bezpiecznych kopii}\label{sec:kryptografia-audyt-backup}

\subsection{Kryptograficzny łańcuch skrótów (Hash Chain)}\label{subsec:hash-chain-teoria}

Zgodnie z wytycznymi NIST SP 800-92 \cite{nist-sp800-92} dzienniki zdarzeń bezpieczeństwa przechowywane na lokalnej maszynie są narażone na celową manipulację (usunięcie niewygodnych wpisów, zmianę znaczników czasu lub podmianę parametrów operacji) w celu zatarcia śladów działania. Aby nadać dziennikowi właściwość wykrywalności manipulacji (\textit{tamper-evidence}), poszczególne rekordy wiąże się w jednokierunkowy łańcuch skrótów kryptograficznych (\textit{hash chain}).

\begin{definition}[Łańcuch skrótów dziennika audytu]\label{def:hash-chain}
Niech $R_1, R_2, \dots, R_n$ oznacza ciąg uporządkowanych rekordów dziennika audytu posiadających ściśle rosnące numery sekwencyjne $\operatorname{seq}(R_i) = i$, a $H \colon \{0,1\}^* \to \{0,1\}^{256}$ oznacza kryptograficzną funkcję skrótu SHA-256. Każdy rekord $R_i$ zawiera pole $\operatorname{prev}_i$ przechowujące skrót poprzedniego rekordu oraz pole $\operatorname{hash}_i$ wyznaczane rekurencyjnie według zależności:
\begin{equation}\label{eq:hash-chain}
    \operatorname{prev}_i = H_{i-1}, \qquad H_i = \operatorname{SHA-256}\!\left(\operatorname{CanonicalJSON}(R_i \setminus \{\operatorname{hash}\})\right), \qquad \text{gdzie } H_0 = 0^{64}.
\end{equation}
gdzie operator $\operatorname{CanonicalJSON}$ oznacza deterministyczną serializację rekordu z leksykograficznym sortowaniem nazw właściwości, niezmiennym formatowaniem znaczników czasu ISO~8601 UTC oraz kodowaniem UTF-8 bez BOM i bez zbędnych białych znaków, co eliminuje zależność skrótu od wewnętrznej kolejności atrybutów w słownikach .NET.
\end{definition}

Z właściwości odporności funkcji SHA-256 na znalezienie kolizji oraz drugiego przeciwobrazu wynika, że jakakolwiek modyfikacja treści pojedynczego rekordu $R_k$ ($1 \le k < n$), usunięcie rekordu ze środka dziennika lub zmiana kolejności wpisów powoduje niezgodność między zapisaną wartością $\operatorname{prev}_{k+1}$ a wyliczonym skrótem $H_k$. Aby zmodyfikować historyczny wpis $R_k$ bez przerwania łańcucha, osoba posiadająca pełny dostęp administracyjny musiałaby przeliczyć i nadpisać wszystkie kolejne skróty $H_k, H_{k+1}, \dots, H_n$ (lub obciąć końcowe wpisy dziennika). Obrona przed całkowitym przepisaniem lub obcięciem końca niekluczowanego łańcucha polega na okresowym utrwalaniu bieżącego nagłówka łańcucha $\langle n, H_n \rangle$ (\textit{head hash}) poza audytowaną maszyną.

\subsection{Atrybucja zmian zewnętrznych w dzienniku zdarzeń Windows}\label{subsec:atrybucja-eventlog}

Gdy zmiana w polityce zapory zostanie wykonana poza skryptem administracyjnym, sam rejestr systemowy przechowuje jedynie nowy stan docelowy, nie zawierając informacji o użytkowniku ani procesie, który dokonał zapisu. Informację tę rejestruje usługa \code{MpsSvc} w kanale dziennika zdarzeń \code{Microsoft-Windows-Windows Firewall With Advanced Security/Firewall}. Po włączeniu podkategorii inspekcji \textit{MPSSVC Rule-Level Policy Change} \cite{ms-audit-mpssvc} system zapisuje w tym kanale zdarzenia opisujące dodanie, modyfikację lub usunięcie reguły wraz z polami XML \code{RuleId}, \code{ModifyingUser} (identyfikator SID sprawcy) oraz \code{ModifyingApplication} (pełna ścieżka pliku wykonywalnego, np.\ \file{mmc.exe}, \file{netsh.exe} lub \file{powershell.exe}). Zestawienie identyfikatorów zdarzeń (różniących się pomiędzy starszymi wydaniami Windows 10 a kompilacją Windows 11 24H2) przedstawiono w tabeli~\ref{tab:event-ids}.

\begin{table}[H]
    \centering
    \vspace{-4pt}
    \caption{Identyfikatory zdarzeń zmian polityki zapory w kanale systemowym Windows}\label{tab:event-ids}
    \tablefont
    \renewcommand{\arraystretch}{0.96}
    \begin{tabularx}{\linewidth}{L{3.6cm} C{3.4cm} Y}
        \toprule
        \hcell{Kategoria zdarzenia} & \hcell{Identyfikatory (Event ID)} & \hcell{Znaczenie operacyjne i pola korelacji} \\
        \midrule
        Dodanie reguły (\code{Added}) & \code{2004}, \code{2097} & Utworzenie nowej reguły w magazynie zapory; zawiera \code{RuleId}, \code{RuleName}, \code{ModifyingUser} oraz \code{ModifyingApplication}. \\
        Modyfikacja reguły (\code{Modified}) & \code{2005}, \code{2098}, \code{2099} & Zmiana parametrów lub stanu aktywności istniejącej reguły o wskazanym \code{RuleId}. \\
        Usunięcie reguły (\code{Deleted}) & \code{2006}, \code{2052} & Usunięcie pojedynczej reguły z magazynu polityki. \\
        Usunięcie wszystkich reguł (\code{AllDeleted}) & \code{2033}, \code{2059} & Masowe wyczyszczenie reguł lub nadpisanie polityki (np.\ podczas importu pliku \file{.wfw} lub resetu zapory). \\
        Zmiana ustawień profilu (\code{Setting}) & \code{2002}, \code{2003} & Zmiana globalnych właściwości profilu sieciowego (\textit{Domain}, \textit{Private}, \textit{Public}). \\
        \bottomrule
    \end{tabularx}
    \vspace{-4pt}
\end{table}

\subsection{Bezpieczeństwo kopii zapasowych w modelu pull i izolacja uprawnień}\label{subsec:backup-pull-teoria}

Tradycyjny model kopii zapasowych typu \textit{push}, w którym stacja robocza może swobodnie zapisywać pliki na serwerze kopii, stwarza poważne ryzyko: przejęcie konta administratora na stacji roboczej pozwala intruzowi połączyć się z serwerem i zniszczyć dotychczasowe kopie zapasowe.

Aby wyeliminować ten wektor ataku, w projekcie zastosowano odwrócony model komunikacji (\textit{server-initiated pull}) połączony z separacją kanałów SSH oraz SFTP \cite{ms-openssh-windows, sftpgo-docs}:
\begin{itemize}[noitemsep, topsep=2pt]
    \item \textbf{Kanał wyzwalania i eksportu (SSH, port 22)} - klucz prywatny stacji Windows nie daje dostępu do powłoki serwera Ubuntu; w pliku \file{authorized_keys} na serwerze jest on powiązany z opcjami \code{from=10.0.2.20,restrict} oraz wymuszonym poleceniem \code{command=/usr/local/bin/fw-server.sh request <user>}. Klient może jedynie poprosić serwer o wykonanie kopii. W odpowiedzi to serwer łączy się zwrotnie ze stacją Windows (klucz eksportu ograniczony w \file{administrators_authorized_keys} do wywołania \code{-Agent Export}) i pobiera strumień konfiguracji.
    \item \textbf{Kanał pobierania kopii przy odtwarzaniu (SFTPGo, port 2022)} - podczas listowania i przywracania kopii stacja Windows łączy się z usługą SFTPGo na porcie 2022. Konto użytkownika w SFTPGo posiada wyłącznie uprawnienia \code{list} oraz \code{download} do wirtualnego katalogu \file{/backups} (brak uprawnień zapisu, nadpisywania i usuwania). Dzięki temu stacja kliencka może odczytać archiwalne kopie i zweryfikować ich skróty SHA-256, lecz w żaden sposób nie może ich usunąć ani zmodyfikować na serwerze.
\end{itemize}
